\documentclass[10pt,letterpaper]{article}
\usepackage{cloudsuite}
\renewcommand{\DocShort}{CH. 04}
\begin{document}
\coverpage{Chapter 4: Fundamental Concepts and Models}{Cornell study notes and applied review | FOUNDATIONS}{Concepts \textbullet{} decisions \textbullet{} verification artifacts}
\sourceStatement
\tableofcontents
\newpage
\section{Learning contract}
\begin{takeaway}
Select a service model and map responsibility and trust boundaries.
\end{takeaway}
\textbf{Chapter scope:} Roles and boundaries, characteristics, delivery and deployment models.\\
\sourceRef{4}{51}
\subsection{Mapped source sections}
\begin{itemize}
\item 4.1 Roles and Boundaries
\item 4.2 Cloud Characteristics
\item 4.3 Cloud Delivery Models
\item 4.4 Cloud Deployment Models
\end{itemize}
\section{Cornell cue and notes matrix}
\small\setlength{\LTleft}{0pt}\setlength{\LTright}{0pt}
\begin{longtable}{>{\RaggedRight\arraybackslash}p{0.29\textwidth}>{\RaggedRight\arraybackslash}p{0.63\textwidth}}
\toprule\textbf{Cue / essential question}&\textbf{Explanatory notes and reasoning}\\\midrule\endfirsthead\toprule\textbf{Cue}&\textbf{Notes}\\\midrule\endhead
\textbf{How do cloud roles differ?} & Provider, consumer, broker, service owner and resource administrator capture distinct operational and contractual viewpoints. One organization may perform multiple roles, but decision rights must be explicit.\\[3pt]\textcolor{petrol}{\emph{Practice:} Build a RACI for service onboarding.}\\[8pt]\midrule
\textbf{What is a trust boundary?} & A trust boundary is a point where expectations, identity, protection or administrative authority change. Network location alone is not a complete security boundary.\\[3pt]\textcolor{petrol}{\emph{Practice:} Annotate boundaries between a client, an API and its data store.}\\[8pt]\midrule
\textbf{What are characteristic cloud capabilities?} & On-demand use, broad access, pooled multitenant resources, elasticity, metered usage and resiliency characterize cloud operations.\\[3pt]\textcolor{petrol}{\emph{Practice:} Propose evidence for each characteristic.}\\[8pt]\midrule
\textbf{How do IaaS, PaaS and SaaS differ?} & Each shifts which portions of infrastructure, runtime or application the provider administers. Responsibility for data, identities and use conditions persists in every model.\\[3pt]\textcolor{petrol}{\emph{Practice:} Fill out a three-model responsibility matrix.}\\[8pt]\midrule
\textbf{How do public, private, hybrid and multicloud relate?} & Deployment arrangements describe placement and ownership relationships, not the depth of service abstraction. Hybrid connects differing environments, while multicloud spans cloud environments.\\[3pt]\textcolor{petrol}{\emph{Practice:} Show one hybrid arrangement and one multicloud arrangement without conflating them.}\\[8pt]\midrule
\textbf{Can models be combined?} & An application can consume IaaS, PaaS and SaaS simultaneously; integration and shared responsibility must be documented across the seams.\\[3pt]\textcolor{petrol}{\emph{Practice:} Identify two cross-model failure or authorization paths.}\\[8pt]\midrule
\end{longtable}\normalsize
\section{Synthesis and architecture connections}
Translate the chapter vocabulary into one explicit service requirement, one boundary, one owner and one measurable test. Later mechanism and architecture chapters provide possible realizations; the requirements remain the acceptance baseline.
\begin{workbox}[title={Applied deliverable}]
Select a service model and map responsibility and trust boundaries.\\[4pt]Create a one-page architecture or operations artifact with: (1) assumptions, (2) a decision and alternatives, (3) control ownership, (4) observable acceptance criteria, and (5) a failed-case test.
\end{workbox}
\section{Active recall}
\begin{enumerate}
\item How do cloud roles differ?
\item What is a trust boundary?
\item What are characteristic cloud capabilities?
\item How do IaaS, PaaS and SaaS differ?
\end{enumerate}
\subsection{Answer-check rubric}
\begin{itemize}
\item A correct response distinguishes the named concepts rather than relying on vendor names.
\item The explanation identifies an operational consequence, boundary or tradeoff.
\item At least one measurement, test or documentary artifact could demonstrate the claim.
\end{itemize}
\section{Blank Cornell capture sheet}
\textbf{Topic / source subsection:}\enspace\rule{0.65\linewidth}{0.4pt}\\[6pt]
\begin{longtable}{p{0.28\textwidth}p{0.64\textwidth}}
\toprule Cues / questions & Notes / evidence \\ \midrule
\rule{0pt}{42pt} & \\ \midrule
\rule{0pt}{42pt} & \\ \midrule
\rule{0pt}{42pt} & \\ \midrule
\rule{0pt}{42pt} & \\ \midrule
\rule{0pt}{42pt} & \\ \midrule
\end{longtable}
\textbf{Summary in three sentences:}\\[16pt]\hrulefill\\[16pt]\hrulefill\\[16pt]\hrulefill

\refs
\end{document}
